\section{RT-X：Cross-Embodiment Data 能否继续扩大 Transfer}

如果单一机器人数据太少，一个自然方向是合并多机构、多机器人、多任务 trajectories。Open X-Embodiment / RT-X 把 action space、observation 和 dataset schema 尽可能统一，再检查一个模型能否从不同 embodiment 受益。这是 scaling law 在 robotics 中更接近现实的试验场。

\subsection{Data mixture 与 action heterogeneity}

Collage 页显示数据来自不同机械臂、camera、workspace 与任务。合并时必须处理 action dimension、coordinate frame、control frequency、gripper semantics 和 missing sensors。简单 concatenate 可能让大 dataset 淹没小 dataset，也可能让 incompatible action labels 产生 negative transfer。

\lecturefigure{slide-37-rtx-data-collage.jpg}{Open X-Embodiment：跨机构、跨机器人数据混合}{官方录像恢复幻灯片 V037；00:56:33}

\lecturefigure{slide-38-rtx-positive-transfer.jpg}{RT-X：跨 embodiment 训练的 positive-transfer 结果}{官方录像恢复幻灯片 V038；00:57:09}

\begin{knowledgebox}{术语消化：robotics scaling 的四个维度}
\begin{tabularx}{\textwidth}{@{}L{0.18\textwidth}X@{}}
\toprule
维度 & 增加的对象与主要风险 \\
\midrule
Model scale & 参数和 compute；风险是 data 不足、latency 与 memory 增大 \\
Task scale & 指令和行为种类；风险是 label/action 冲突与 long-tail failure \\
Environment scale & 场景、物体和背景；风险是 observation shift 与新 affordance \\
Embodiment scale & 机器人形态和 action space；风险是坐标、频率与控制语义不兼容 \\
\bottomrule
\end{tabularx}
\end{knowledgebox}

\teachervoice{Fei Xia 对 RT-X 的总结很克制：机器人领域看到了 positive transfer 的“signs of life”，但语言模型社区已经把 transfer 当作常态。换言之，cross-embodiment scaling 是重要信号，却还没有到“更多数据必然更好”的成熟阶段。}

\subsection{Part 1 的结论}

第一部分最终收束为三点：model consolidation 让 high-level reasoning 与 low-level control 更接近同一系统；joint scaling 测试 diverse data 是否产生 transfer；positive transfer 是评估 scaling 是否成立的关键。接下来课程不再继续扩大一个 end-to-end policy，而是换一种接口，把 LLM 的输出改成 reward code。

\lecturefigure{slide-39-part1-summary.jpg}{Part 1：model consolidation、joint scaling 与 positive transfer}{官方录像恢复幻灯片 V039；00:57:51}

\subsection{本章小结}

RT-X 把数据稀缺问题转化为跨 embodiment 数据标准化与 transfer 问题。它扩展了 RT-2 的路线，但 action heterogeneity 和 dataset imbalance 让“统一模型”更难；这也解释为什么第二条路线选择保留传统 controller。

